\section{Foundation Agent: de lenguaje a acción}

\subsection{Aplicación de los modelos base en los Agentes}

Jim Fan propone el concepto de \textbf{Foundation Agent}: utilizar modelos base preentrenados a gran escala para construir Agentes de propósito general:

\begin{itemize}
    \item \textbf{Voyager}: un Agente de exploración de mundo abierto en Minecraft que usa GPT-4 como núcleo de razonamiento, descubre habilidades automáticamente, escribe código y lo almacena en una biblioteca de habilidades, logrando aprendizaje continuo
    \item \textbf{Eureka}: usa un LLM para generar y optimizar automáticamente la reward function, entrenando manipulación diestra con manos en simulación física
    \item \textbf{MineDojo}: una base de conocimiento de Minecraft a gran escala y un banco de pruebas de mundo abierto
\end{itemize}

\begin{knowledgebox}{La biblioteca de habilidades de Voyager}
Cada habilidad que Voyager aprende (como fabricar una espada o construir una casa) se almacena en forma de código ejecutable. Las tareas posteriores pueden invocar directamente las habilidades ya existentes, sin necesidad de aprender desde cero; este es un excelente ejemplo de memoria a largo plazo y aprendizaje continuo en los Agentes.
\end{knowledgebox}

\subsection{Resumen de la sección}

Mediante las capacidades de razonamiento y generación de código de los LLM, Foundation Agent logra exploración de mundo abierto y aprendizaje continuo de habilidades en mundos virtuales.

